\section{The argument system}
\label{sec:argument}

This section states what is proved and how, independently of how any component
is built. A \emph{shard argument} reduces every claim about one shard's
execution to a single evaluation of a single committed polynomial, and one
opening authenticates it. A \emph{block argument} composes shard arguments by
recursion into one proof of the relation $\mathcal R$ of
Section~\ref{sec:statement}. Sections~\ref{sec:air} to~\ref{sec:recursion}
instantiate each component named here, and Section~\ref{sec:verification} asks
whether the instantiation admits only executions of the advertised machine.

\subsection{Parameters}
\label{sec:argument-params}

The protocol is stated over the following parameters, and is symbolic in all of
them: a different choice gives a different instantiation of the same protocol.
Numeric values are fixed where the component is built.
Table~\ref{tab:whir-schedule} fixes the proximity-test schedule of the
checked-in core and compress profiles, \S\ref{sec:pcs-geometry} the stacking
height and the row cube, \S\ref{sec:prelim} the permutation schedule,
\S\ref{sec:isa} the clock fence and \S\ref{sec:recursion} the recursion
arity. The measured baseline of Section~\ref{sec:performance} differs from the
checked-in configuration in two of these parameters, and one of them is a
security parameter: it sets the lookup grinding to zero, and it carries the
degree-seven $\Poseidon{}$ round count rather than the corrected degree-three
schedule (\S\ref{sec:prelim}).

\begin{itemize}
  \item The base field $\F$ and the extension $\EF$ from which challenges are
  drawn. The extension degree bounds every soundness term of the form
  $\deg/|\EF|$.
  \item The permutation used for the transcript and for the commitment trees,
  with its round counts, and the arity of those trees.
  \item The stacking height $H=2^{\ell}$, the row cube $2^{n}$, and the bucket
  rule that maps a round's real area to its committed area
  (\S\ref{sec:pcs-geometry}). These fix the committed geometry, and with it
  the set of recursion programs a prover can produce.
  \item For the proximity test: the initial rate $\rho$, the per-round folding
  factors, the query counts, the number of out-of-domain samples, the grinding
  bits spent on each query phase and on the bus challenges, and the regime in
  which the query soundness is counted.
  \item For recursion: the arity of a compose node and the height of the tree of
  admissible recursion keys (\S\ref{sec:recursion}).
  \item For the machine: the configured shard size, the clock fence and the
  per-unit height budgets, whichever of which binds first closing a shard
  (\S\ref{sec:isa}).
  \item The guest-accelerator set, which fixes the precompile units a shard may
  contain and therefore the unit set itself (\S\ref{sec:ipcore-config}).
  \item For the cross-shard digest: the curve and its extension, the bound
  $N_{\mathrm{int}}$ on the interactions a block proof may contain, and the
  per-interaction multiplicity bound $B$ (\S\ref{sec:air-digest-security}).
\end{itemize}

Two properties of the instantiation are worth stating with the parameters
rather than leaving to a discussion. Query soundness is counted in the
unique-decoding regime, where no conjecture is needed, rather than at the
Johnson bound or towards capacity (\S\ref{sec:iopp-soundness}). And no
blinding is applied anywhere: the protocol is succinct, not hiding, and a proof
reveals what its committed trace reveals.

\subsection{Syntax and components}
\label{sec:argument-components}

We state the shard argument and each primitive it uses as a tuple of
algorithms, so that the protocol of \S\ref{sec:argument-shard} can be read
against interfaces rather than against any one construction. Throughout,
$\mathsf{pp}$ are the public parameters of \S\ref{sec:argument-params} and
$\tau$ is the transcript the Fiat--Shamir transform maintains.

A \emph{shard argument} is a tuple $\mathsf{SA} = (\mathsf{Setup},
\mathsf{Prove}, \mathsf{Verify})$:

\begin{description}
  \item[$\mathsf{Setup}(\mathsf{pp}, P) \to (\mathsf{pk}, \mathsf{vk})$:] on
  input a program image $P$, commits to the program ROM and the preprocessed
  tables and outputs a proving key and a verifying key.
  \item[$\mathsf{Prove}(\mathsf{pk}, T) \to \pi$:] on input the execution
  record $T$ of one shard, outputs a shard proof.
  \item[$\mathsf{Verify}(\mathsf{vk}, \mathsf{pv}, \pi) \to b$:] on input the
  shard's public values, outputs $b \in \{0,1\}$.
\end{description}

It is built from four primitives. A \emph{unit} is a table of columns over
$\F$ with a constraint system and a set of bus interactions
(Definitions~\ref{def:chip} and~\ref{def:bus}), and a shard is a set of units
with a shared public-value row; units are data rather than algorithms, so they
carry no syntax here. The other three do.

A \emph{bus argument} $\mathsf{BA} = (\mathsf{Prove}, \mathsf{Verify})$
certifies that the multiset of tuples sent on each bus equals the multiset
received:

\begin{description}
  \item[$\mathsf{BA.Prove}(\tau, \{U_i\}) \to (\pi_{\mathrm{bus}}, \bm z, \{v_{i,j}\})$:]
  outputs a proof together with the point $\bm z$ and the claimed evaluations
  $v_{i,j}$ of column $j$ of unit $i$ at $\bm z$.
  \item[$\mathsf{BA.Verify}(\tau, \pi_{\mathrm{bus}}) \to (\bm z, \{v_{i,j}\}) \lor \bot$:]
  accepts and outputs the point and the claims it has reduced to, or rejects.
\end{description}

A \emph{constraint argument} $\mathsf{CA} = (\mathsf{Prove}, \mathsf{Verify})$
has the same shape and certifies instead that every unit's constraints vanish
on its real rows, reducing that claim to evaluations at the same point.

A \emph{jagged commitment} $\mathsf{JC} = (\mathsf{Commit}, \mathsf{Reduce},
\mathsf{Open}, \mathsf{VerifyOpen})$ carries columns of differing heights:

\begin{description}
  \item[$\mathsf{JC.Commit}(\mathsf{pp}, (c_1,\dots,c_M)) \to (\mathsf{com}, \mathsf{st})$:]
  on input $M$ columns of arbitrary heights, places them into the committed
  geometry and outputs a commitment and prover state.
  \item[$\mathsf{JC.Reduce}(\mathsf{st}, \{(i, \bm z, v_i)\}) \to (\pi_{\mathrm{jag}}, \bm z^{*}, v^{*})$:]
  reduces any set of column-evaluation claims to a single evaluation
  $v^{*}$ of one dense multilinear polynomial at one point $\bm z^{*}$.
  \item[$\mathsf{JC.Open}(\mathsf{st}, \bm z^{*}) \to \pi_{\mathrm{open}}$] and
  \item[$\mathsf{JC.VerifyOpen}(\mathsf{com}, \bm z^{*}, v^{*}, \pi_{\mathrm{open}}) \to b$:]
  authenticate that single evaluation against the commitment.
\end{description}

Only the last interface constrains the others: because every claim is reduced
to one evaluation of one polynomial, the number of units, their widths and
their heights are invisible to the opening, and a unit may be added, widened or
removed without changing the layer above it.

\subsection{The shard argument}
\label{sec:argument-shard}

Protocol~\ref{prot:shard} states the reduction.
Each step ends with claims of the same kind it began with, evaluations of
committed columns, so the steps compose without the verifier ever reading a
column.

\begin{protocol}[Shard argument]
\label{prot:shard}
\textbf{Statement:} a verifying key committing to the program and the
preprocessed tables, and the public values of one shard.\\
\textbf{Witness:} the shard's execution record.
\begin{enumerate}
  \item \textbf{Commit.} The prover arranges each unit's events into rows,
  places the columns of each commitment round into one dense vector and
  commits to it. Round order is fixed and the per-unit dimensions enter the
  transcript, so the geometry is public before any challenge is drawn
  (Section~\ref{sec:pcs}).
  \item \textbf{Buses.} The verifier draws the bus challenges. The prover runs
  the bus argument, which reduces every send and receive of every unit to
  column evaluations at one point (Section~\ref{sec:air-logup}).
  \item \textbf{Constraints.} The verifier draws the constraint challenges. The
  prover runs the constraint argument on the same point, so each column is
  opened once for both arguments (Section~\ref{sec:air-zerocheck}).
  \item \textbf{Jagged reduction.} The column claims of the two previous steps,
  together with explicit zero claims for the gaps in the committed geometry,
  are batched and reduced to one evaluation of the dense polynomial
  (Section~\ref{sec:pcs-jagged}).
  \item \textbf{Assist.} The verifier does not materialise the weight table the
  previous step reduced against. A second sumcheck over the public column
  boundaries certifies it, and any point-extension coordinates are sampled
  here, in an order that is itself part of the protocol
  (Section~\ref{sec:pcs-jagged-assist}).
  \item \textbf{Opening.} One batched opening authenticates that evaluation
  against every commitment of step~1 (Section~\ref{sec:iopp}).
\end{enumerate}
\textbf{Output:} a shard proof whose public values carry the entry and exit
machine state, the memory-address cursors, the digest of committed output and
the exit code that the relation of Section~\ref{sec:statement} is stated over,
and a digest of the interactions this shard leaves unbalanced.
\end{protocol}

\paragraph{From shards to a block.} A shard argument proves one shard in
isolation. Three things tie the shards of one execution together, and each is
discharged where it arises rather than by the opening: the machine state
chains through the public values, the memory-address cursors tile, and the
interactions that cross a shard boundary accumulate into a digest that must
cancel over the whole execution (Section~\ref{sec:air-global}). A recursion
layer then verifies shard proofs inside the argument system itself and merges
them, so that the verifier of a block reads one proof rather than one per
shard (Section~\ref{sec:recursion}).

\paragraph{What is proved.} Two statements carry the construction, and both are
conditional in ways the paper makes explicit rather than hiding.
Theorem~\ref{thm:composition} bounds a cheating prover for the whole tree by
the number of random-oracle queries times the sum of the per-node errors, plus
the digest term, under a per-node round-by-round premise and a hypothesis about
composed extractors. Theorem~\ref{thm:unique} shows that two satisfying traces for the same
program and public values agree up to permutation, given unit determinism and
the real-row premises that Section~\ref{sec:verification} reports coverage
for; identifying that trace with an execution of the machine needs the
conformance evidence of the same section.
Neither statement is about a particular unit set, a particular field or a
particular proximity test; those are the instantiation, fixed where each
component is built.
